\subsection{Long-tail recognition and classifier rebalancing}

Long-tail recognition has developed several ways to separate representation
learning from classifier fitting. These include class-balanced classifier
retraining, decoupled representation/classifier procedures, score or logit
adjustment, and corrections based on classifier norms or class frequencies
\cite{kang2020decoupling,alshammari2022weight,menon2021logit}. Such methods establish that a classifier head can be a useful
intervention point after or alongside representation learning. They also make
clear that balanced head retraining is not, by itself, a new contribution.

These approaches differ in their optimization goals and deployment intent, but
they share the premise that class-frequency effects need not be addressed only
through the representation. They provide important controls for interpreting a
rare-class improvement: a change after head adaptation need not imply that the
feature extractor has changed. At the same time, a nominal tail-class gain does
not specify where a mistaken ordinal decision lies relative to its true class.
This gap motivates the localization-focused endpoint used here.

Our analysis uses this literature as motivation, not as an algorithmic target.
The direction-only intervention is not proposed as a replacement for balanced
retraining or as a general long-tail classifier. It is a constrained probe that
holds the representation, original row norms, and biases fixed. This lets the
analysis ask whether rare-end ordinal localization can change with the
direction component of the head. The primary outcome is an ordinal localization
endpoint rather than a generic tail-class accuracy measure.

The constraint is also important for the paper's claim boundary. Full-head
retraining can change directions, norms, and biases simultaneously, making it
hard to attribute a localization change to one head component. The present
comparison is not intended to rank all rebalancing strategies. It tests whether
a direction change remains associated with rare-end movement when these other
head quantities are fixed.

\subsection{Classifier and representation geometry under imbalance}

Prior work has examined angular classifier structure, classifier weight norms,
long-tail representation geometry, and alignment between features and
classifier weights
\cite{c2am2022angular,zhu2022balanced,yi2025geometry,wang2026alignment}. These studies establish that feature and head
geometry can matter under imbalance. The present work does not claim that
classifier direction, norm effects, representation distortion, or
feature/classifier misalignment are new observations.

Geometry can be used at several levels of analysis. It can describe the
separation of class features, the orientation of classifier weights, or the
relation between a feature and multiple class prototypes. Our use of
train-derived centroids is deliberately limited to a diagnostic role. It does
not turn nearest-centroid routing into a ground-truth label for representation
failure, and it does not imply that a sample without exact recovery lacks useful
endpoint information.

Our focus is more specific. We study rare upper endpoints in an ordered label
space and measure their movement toward adjacent interior classes. The
fixed-norm, fixed-bias direction-only comparison tests a head-actionable part
of this localization under frozen representations. We then ask whether an
endpoint-versus-adjacent centroid margin adds information about final
post-adaptation localization beyond original-head state and generic centroid
separation. The result is deliberately reported with its boundary: the geometry
diagnostic is strong in Solar but partial or unsupported in the two
RetinaMNIST settings. It is therefore not presented as a universal
geometry-conditioned recovery mechanism.

This distinction is central to the interpretation of the negative results.
The geometry analysis is not a post-hoc criterion for selecting successful
samples or a substitute for the cross-seed H1 replication. It is a fixed
diagnostic layered on top of a replicated intervention result. The strong Solar
outcomes and mixed RetinaMNIST outcomes therefore delimit the claim rather than
forming evidence for a general geometric law.

\subsection{Ordinal classification under imbalance}

Ordinal classification methods account for label order through structured
losses, thresholds, costs, ranking objectives, or representation constraints.
Imbalance-aware variants also use weighting, resampling, or other mechanisms to
improve minority-class predictive performance
\cite{zhu2019smor,lazaro2023bayesian}. These approaches
address important aspects of ordered prediction, but their usual objective is
overall classification or ordinal-error improvement. They need not distinguish
whether errors of a rare upper endpoint move one step inward, several steps
inward, or remain centered below the endpoint in the predictive distribution.

The ordered setting also makes endpoint asymmetry relevant. A lower and upper
endpoint can have different local neighborhoods and different directions of
possible displacement. The paper focuses on the rare upper endpoint because its
inward movement is directly measurable along the observed ordered scale. This
focus does not imply that lower-end behavior is unimportant or that all ordinal
imbalance problems share the same asymmetry.

This paper instead treats rare-end inward localization as the empirical object
of analysis. It uses routing, predictive location, shrinkage, and rare-end
ordinal error to describe the phenomenon, then separates descriptive
representation/head diagnostics from a replicated head intervention. The aim
is not to claim that existing ordinal imbalance methods omit all endpoint
analysis, but to provide a controlled characterization of this particular
failure mode.

\subsection{Position of this work}

The paper is a mechanism and empirical analysis study. It does not propose a
new rebalancing algorithm, claim that classifier direction is novel, or assert
that geometry universally explains rare-end recovery. Its contribution is to
characterize an ordinal endpoint-localization failure, show a reproducible
fixed-norm, fixed-bias direction-only response in the tested RetinaMNIST and
Solar settings, and report the setting-dependent limits of the associated
geometry diagnostic.
